% =====================================================================
% Determinacy-Field Mechanics, Paper 3 -- SECTION SKELETON (v0.1, 2026-08-16)
% Theme: "Reality calibration" -- adopted 2026-08-16 (原仮定者裁定, 第121-122便).
%   The emergence-across-scales draft that previously held this slot moved,
%   content-unchanged, to Paper 4 (dfm-paper4.tex).
% Status: chapter outline (章立て素案) with the measured anchor numbers from
%   the landed harnesses. All prose is bullet-level; nothing here is final.
% Pillars (chronological by wave):
%   P1  43''/century reproduced by direct dynamics (waves 113-115:
%       real-c convention, stateCarry:"double"; ratio 1.002).
%   P2  kF1 stability is unit-system dependent -- the loop gain is built
%       from the DIMENSIONAL weight w=m/d (wave 119).
%   P3  kF1 drag precession as a calibration mechanism: the lunar apsidal
%       cycle 8.85 yr reproduced by a two-body system via motion dragging,
%       with D0 as the calibration knob (wave 120).
%   P4  The joint correction (D0=0.006, q=5): one pair satisfies both the
%       Earth-Moon and Mercury systems at once; the spin term is suppressed
%       geometrically by q, and the wave-121 seven-order D0 split dissolves
%       (wave 122, tests/exp-kf1c.mjs).
% Honest-calibration philosophy (runs through every section):
%   calibrated match =/= mechanism identification; every failure that led
%   to a convention (kFrame=0 samples, softening retrograde precession,
%   window dependence) is reported as part of the result.
% Companion papers: Paper 1 (foundations), Paper 2 (box universe),
%   Paper 4 (emergence across scales, dfm-paper4.tex).
% Measurement harnesses of record: tests/exp-obscal.mjs, exp-kf1.mjs,
%   exp-kf1b.mjs, exp-kf1c.mjs; app samples: mercuryReal (comet emoji
%   in-app), earthMoonReal, earthMoonRealKF1, mercuryRealKF1,
%   saturnRingReal (docs/PHYSICS.md sec. 5, v1.34).
% =====================================================================
\documentclass[aps,reprint,amsmath,amssymb,nofootinbib,superscriptaddress,floatfix]{revtex4-2}

\usepackage{graphicx}
\usepackage{hyperref}
\usepackage{booktabs}

\begin{document}

\title{Reality Calibration of a Machian Toy Universe:\\
From 43$''$/century to the Lunar Apsidal Cycle\\
by Honest Calibration, Not Mechanism Claims}

% TODO(author): author block as in Papers 1--2.
\author{...}
\date{2026-08-16 (skeleton v0.1)}

\begin{abstract}
% --- Planned claims (each tied to a harness measurement) ---
% 1. A single scale-conversion convention (per-sample exponents L/T/M,
%    real G = 6.674, real c0 = 3x10^4, Kt = c0^2/G "correspondence lock")
%    lets real solar-system data be placed in the toy universe unchanged.
% 2. With the double-precision state carry, the 1PN geodesic mode
%    reproduces Mercury's 42.98''/century directly in the dynamics
%    (ratio 1.002) -- not as a formula overlay (P1).
% 3. The kFrame=1 sector calibrates: motion dragging reproduces the lunar
%    apsidal cycle (8.85 yr) in a PURE TWO-BODY system (P3), and the joint
%    correction (D0=0.006, q=5) makes the same pair of constants hold for
%    the Earth-Moon and Mercury systems simultaneously (P4).
% 4. Throughout: what a calibrated match does and does not license.
%    The real lunar apsidal driver is solar perturbation; DFM matches the
%    observable via a different mechanism. We claim calibration agreement,
%    never mechanism identification.
% --- Honest negatives (fixed list, planned Sec. VIII) ---
%    D0 remains dimensional (P2); q=5 trades away the Lense-Thirring
%    distance law of the spin term; window dependence of the apsidal
%    ratio; the dragged body's own spin is provably irrelevant
%    (self-excluded from u) -- a hypothesis tested and honestly retired.
TODO abstract.
\end{abstract}

\maketitle

% =====================================================================
\section{Introduction}
% - Papers 1/2 recap in two sentences; this paper asks: how far can a
%   relational toy universe be pushed TOWARD REAL NUMBERS before it breaks,
%   and what must be said honestly at each step?
% - Define "reality calibration" vs "mechanism identification" up front.
% - Chronology-as-method: each pillar P1-P4 arose from a failure of the
%   previous convention; the failures are data.
TODO.

% =====================================================================
\section{Scale-conversion conventions}
% - Per-sample scale exponents {L,T,M}; display-only, physics-invariant.
% - Real-G convention (G_sim* per tier), real-c convention (c0=3x10^4),
%   correspondence lock Kt=c0^2/G (physLock), massFloor=1e-6.
% - eC = L-T for real-c presets; the 9-row scale base table.
% - Table: sample-by-sample exponents and the SI values they pin.
TODO.

% =====================================================================
\section{P1: 43$''$/century by direct dynamics}
% - stateCarry:"double" (wave 114): sub-ulp velocity carry; measured
%   42.98''/century, ratio 1.002 vs GR's 42.98.
% - The invariant ladder (wave 113) and softening retrograde precession
%   (wave 120: eps=2 gave -1.2e-4 rad/orbit = 24x 1PN -- the honest reason
%   the raw HUD needed a subtraction method).
TODO.

% =====================================================================
\section{P2: Stability is unit-system dependent}
% - Loop gain g=[w_sat/(D0+w_sat)][w_M/(D0+w_M)] with w=m/d DIMENSIONAL;
%   the same configuration collapses or survives depending on the unit
%   system (wave 119, exp-kf1.mjs).
% - This is the root of every D0 subtlety below -- state it early.
TODO.

% =====================================================================
\section{P3: The lunar apsidal cycle from motion dragging}
% - kF1 drag precession is motion-drag origin (spin-independent,
%   lambdaPN-independent), proportional to chi = w/(D0+w).
% - Joint fit D0=0.006 + initial speed -0.32%: sidereal month 27.33 d
%   (+0.03-0.04%) AND apsidal cycle 8.85 yr simultaneously, two bodies only.
% - Honest note as a section, not a footnote: the real driver is solar
%   perturbation; a pure two-body match is a DIFFERENT mechanism.
TODO.

% =====================================================================
\section{P4: The joint correction (D0=0.006, q=5)}
% - Mercury's kF1 retrograde precession is the SPIN term
%   omega(d) = s (R/(R+d))^q; suppression factor (R/(R+a))^(q-3);
%   derived q* = 3 + ln(|A3|/eps)/ln((R+a)/R) = 4.59 -> adopt q=5.
% - Measured: drag -5.3e-8 rad/orbit (1/10 of 1PN) at q=5; Earth-Moon
%   ratio purifies 0.988 -> 0.999 (the Earth-spin contamination leaves).
% - ONE pair (D0, q) satisfies both systems (exp-kf1c joint verdict).
% - The wave-121 interim answer (D0=1e4, seven orders from the Moon's
%   0.006) is reported as the failure that motivated this section.
% - Dragged-side spin: x10^4 changes nothing (self-excluded from u) --
%   the hypothesis was tested and is retired honestly.
TODO.

% =====================================================================
\section{Ring-system scope (Saturn)}
% - saturnRingReal at e6 with six major moons: Titan 15.949 d vs 15.945 d,
%   C-ring 5.78 h vs 5.8 h, 22.5 km/s -- Kepler-level calibration only;
%   Lense-Thirring scope stays excluded (kFrame=0 sample).
TODO.

% =====================================================================
\section{What calibration licenses -- and what it does not}
% - The philosophy section: calibrated match vs mechanism identification;
%   D0's dimensionality imports a system scale; q=5 trades the LT distance
%   law; window dependence quantified.
TODO.

% =====================================================================
\section{Negative claims}
% - Mirrors Papers 1/2/4: no claim about real gravitation, no 3D
%   extrapolation, no solar-perturbation physics, etc. (fixed list).
TODO.

% =====================================================================
\section{Reproducibility}
% - Harnesses exp-obscal/kf1/kf1b/kf1c, result JSONs, app sample IDs,
%   QA gates (wave120-122.ui), commit/tag/DOI at submission time.
TODO.

\begin{thebibliography}{9}
% TODO: Papers 1, 2, 4 + standard references (CODATA, IAU) as needed.
\end{thebibliography}

\end{document}
